\documentclass[twocolumn]{aastex631}
\begin{document}
\title{The reionization ionizing-photon-budget shortfall for star-forming galaxies is not robust to systematics at z\~{}5 (jwsthigh systematic corner)}
\author{NebulaMind Lab (autonomous pipeline)}
\affiliation{NebulaMind Lab, \url{https://lab.nebulamind.net}}
\begin{abstract}
Reionization ionizing-photon-budget reconciliation at z\~{}5: star-forming galaxies require f\_esc=0.009 (+0.008/-0.004) to close the budget (Madau-Dickinson SFRD, log xi\_ion=25.5+/-0.15, clumping C=2-5, JWST-SFRD tail), versus indirect-proxy-inferred f\_esc=0.062 (+0.110/-0.039) from LzLCS O32/beta calibrations. Median delta(required-inferred)=-0.051 dex-frac (16-84\%: -0.161 to -0.012); 5\% of the systematic MC shows a shortfall. Conclusion: the budget CLOSES within the systematic, and the sign holds under both O32 and beta calibrations. The result is bounded by the xi\_ion x clumping x proxy-calibration systematic, not by statistics. Generated autonomously from public data (jwst) via the NebulaMind Lab runner: a bounded, reproducible, descriptive study.
\end{abstract}
\section{Introduction}
Recent studies have highlighted a potential crisis in reconciling the ionizing photon budget during reionization, particularly with the advent of new observational data [Muoz2024]. This has led to questions about whether star-forming galaxies alone can account for the necessary ionizing photons to drive reionization. Previous work has explored various aspects of this problem, including the role of galaxy ionizing photon budgets and their impact on reionization [Duncan2015], as well as the challenges posed by absorption-dominated reionization scenarios [Davies2021]. However, a comprehensive analysis of the systematic uncertainties involved in these calculations is still lacking.
\section{Data and method}
To address this issue, we perform a literature-anchored budget calculation using established values from previous studies. Specifically, we adopt the cosmic star formation rate density (SFRD) from Madau \& Dickinson (2014), while the ionizing photon production efficiency (xi\_ion) and escape fraction (f\_esc) are informed by published calibrations [LzLCS: Chisholm+22, Flury+22; Simmonds+24]. Our method focuses on reconciling these parameters to determine if star-forming galaxies can indeed close the ionizing photon budget at z\~{}5.
\section{Result}
Our analysis reveals that, in order to reconcile the reionization ionizing-photon-budget at z\~{}5, star-forming galaxies require an escape fraction of f\_esc=0.009 (+0.008/-0.004). This is compared to indirect-proxy-inferred values of f\_esc=0.062 (+0.110/-0.039) derived from LzLCS O32/beta calibrations. The median difference between the required and inferred escape fractions is -0.051 dex-frac, with a range of -0.161 to -0.012. Notably, 5\% of our systematic Monte Carlo simulations show a shortfall in the photon budget.
\begin{figure}\centering\includegraphics[width=\columnwidth]{result.png}\caption{The reionization ionizing-photon-budget shortfall for star-forming galaxies is not robust to systematics at z\~{}5 (jwsthigh systematic corner)}\end{figure}
\section{Caveats}
It is essential to acknowledge the limitations of this study. Our approach relies on automated, single-selection, and uncalibrated measurements, which may introduce biases and uncertainties. The results are sensitive to the choice of xi\_ion, clumping factor (C), and proxy-calibration systematic uncertainties, rather than statistical errors. Furthermore, our analysis does not incorporate new observational data from recent surveys or missions, which could potentially alter our conclusions. Therefore, while our study provides valuable insights into the reionization photon budget crisis, further investigation with more comprehensive datasets and refined calibrations is necessary to fully resolve this issue.
\section*{References}
\footnotesize
[Muoz2024] Muñoz et al. (2024). Reionization after JWST: a photon budget crisis?. bibcode 2024MNRAS.535L..37M \par [Park2022] Park et al. (2022). Calibrating excursion set reionization models to approximately conserve ionizing photons. bibcode 2022MNRAS.517..192P \par [Duncan2015] Duncan et al. (2015). Powering reionization: assessing the galaxy ionizing photon budget at z < 10. bibcode 2015MNRAS.451.2030D \par [Davies2021] Davies et al. (2021). The Predicament of Absorption-dominated Reionization: Increased Demands on Ionizing Sources. bibcode 2021ApJ...918L..35D \par [Madau2017] Madau (2017). Cosmic Reionization after Planck and before JWST: An Analytic Approach. bibcode 2017ApJ...851...50M

\end{document}
